\section{Localization as a certified setoid}

\subsection{Fractions and equality witnesses}

Let $A$ be a certified commutative ring and fix $f\in A$.  A code for an element of $A_f$ is a pair
$(a,n)$, written $a/f^n$.  The pair is valid when $a$ is a valid $A$-code.  Equality is not defined by
normalizing fractions.  A witness from $a/f^n$ to $b/f^m$ is
\begin{enumerate}[label=(\alph*)]
\item a natural number $k$, and
\item an $A$-equality witness for
\begin{equation}\label{eq:loc-equality}
             f^k\bigl(f^m a-f^n b\bigr)=0.
\end{equation}
\end{enumerate}
This is the standard equality relation for localization at
$\{1,f,f^2,\ldots\}$, made proof-relevant.

\begin{theorem}[Localization witness calculus]\label{thm:localization-calculus}
The relation \eqref{eq:loc-equality} is an equivalence relation.  Addition and multiplication
\[
 \frac a{f^n}+\frac b{f^m}
   =\frac{f^m a+f^n b}{f^{n+m}},\qquad
 \frac a{f^n}\frac b{f^m}=\frac{ab}{f^{n+m}}
\]
are certified congruent operations.  All witness transformers are primitive recursive.
\end{theorem}

\begin{proof}
For reflexivity choose $k=0$ and the certificate for $f^n a-f^n a=0$.  Symmetry uses the same
$k$ and negates the base equality.

For transitivity, suppose
\begin{align*}
 f^k(f^m a-f^n b)&=0,\\
 f^\ell(f^r b-f^m c)&=0.
\end{align*}
Multiply the first identity by $f^{\ell+r}$ and the second by $f^{k+n}$, then add.  The middle
terms cancel and give
\[
 f^{k+\ell+m}(f^r a-f^n c)=0.
\]
Thus a valid killing exponent for the composite witness is $k+\ell+m$, where $m$ is the
denominator exponent of the middle fraction.  The new base-ring witness is constructed from the old
ones using multiplication congruence, addition, and the distributive-law certificates.

For addition and multiplication, expand the cross-multiplied difference of the outputs.  It is a
finite sum of multiples of the two input differences.  If input witnesses kill those differences by
$f^k$ and $f^\ell$, then multiplying by a sufficiently large explicit power, for example
$f^{k+\ell+n+m+n'+m'}$, kills the output difference.  A sharper exponent is easy but unnecessary.
The expansion and witness combination are fixed polynomial identities, hence primitive recursive.
\end{proof}

The theorem remains valid with zero divisors.  The killing exponent is essential precisely because
cross multiplication need not be cancellative.

\subsection{Localization of modules}

For an $A$-module $M$, represent an element of $M_f$ by $x/f^n$ and witness equality by
\[
             f^k(f^m x-f^n y)=0_M.
\]
The preceding proof uses only additivity and scalar multiplication, so it carries over verbatim.
The restriction map $M\to M_f$ sends $x$ to $x/f^0$.  Equality of two restricted elements is
witnessed by a local annihilation equation $f^k(x-y)=0$.

If $g\in A$, restriction from $M_f$ to $M_{fg}$ sends
\[
             \frac{x}{f^n}\longmapsto \frac{g^n x}{(fg)^n}.
\]
This formula is preferable to an abstract universal property in an object-theory implementation:
it is a fixed primitive-recursive map, and its functoriality is a finite polynomial identity.

\subsection{Universal property in certificate form}

Suppose $\phi:A\to B$ is a certified ring map and $u$ is a supplied inverse witness for $\phi(f)$,
that is, a code $v\in B$ with certificates
\[
          \phi(f)v=1,
          \qquad v\phi(f)=1.
\]
Then
\[
       \widetilde\phi(a/f^n)=\phi(a)v^n
\]
defines a certified map $A_f\to B$.  If \eqref{eq:loc-equality} holds, applying $\phi$ and multiplying
by the appropriate power of $v$ proves equality of the images.  Uniqueness is checked on a supplied
fraction by the equation
\[
 a/f^n=(a/1)(1/f)^n.
\]
Thus the universal property can be used in iPRA as a family of finite witness transformations,
without quantifying over a completed class of all ring maps.

\subsection{Localizing a presented quotient}

Let $A=R[X]/I$ be represented by ideal-combination witnesses.  A localization equality witness for
$p/f^n$ and $q/f^m$ consists of a killing exponent $k$ and multipliers $h_i$ satisfying
\[
 f^k(f^m p-f^n q)=\sum_i h_i g_i
\]
in the polynomial ring.  This is a single coefficientwise certificate.  Hence quotient formation and
localization compose without choosing normal forms.

The same presentation can be flattened algebraically:
\[
        (R[X]/I)_f \cong R[X,T]/(I,\,Tf-1).
\]
The map from a fraction sends $p/f^n$ to $pT^n$; the inverse sends $T$ to $1/f$.  Both composites have
explicit relation witnesses.  The fraction presentation is usually smaller for gluing, while the
polynomial presentation is convenient for elimination and finite-presentation libraries.

\subsection{Primitive-recursive implementation and induction audit}

The code operations needed here are bounded exponentiation, tuple access, polynomial addition and
multiplication, and witness-vector combination.  Exponentiation is primitive recursion on $n$ with
the quantifier-free invariant that the accumulator equals the requested finite power.  Polynomial
operations are bounded folds over finite lists.  No induction formula asserts that two localization
classes are equal; it asserts only that a concrete, checked equality witness has been transformed
correctly.
